% !TEX TS-program = XeLaTeX
% use the following command:
% all document files must be coded in UTF-8
\documentclass[portuguese]{textolivre}
% build HTML with: make4ht -e build.lua -c textolivre.cfg -x -u article "fn-in,svg,pic-align"

\journalname{Texto Livre}
\thevolume{19}
%\thenumber{1} % old template
\theyear{2026}
\receiveddate{\DTMdisplaydate{2026}{1}{28}{-1}} % YYYY MM DD
\accepteddate{\DTMdisplaydate{2026}{4}{24}{-1}}
\publisheddate{\DTMdisplaydate{2026}{9}{11}{-1}}
\corrauthor{Cleomar Poletto}
\articledoi{10.1590/1983-3652.2026.64114}
%\articleid{NNNN} % if the article ID is not the last 5 numbers of its DOI, provide it using \articleid{} commmand 
% list of available sesscions in the journal: articles, dossier, reports, essays, reviews, interviews, editorial
\articlesessionname{articles}
\runningauthor{Poletto} 
%\editorname{Leonardo Araújo} % old template
\sectioneditorname{Daniervelin Pereira~\orcid{0000-0003-1861-3609}}
\layouteditorname{Rayany Chaves~\orcid{0009-0007-0028-484X}}

\title{Narradores tradicionais do Vale do Jequitinhonha – (MG): guardiões da memória e da ancestralidade: interfaces da literatura oral, variação linguística e tecnologias digitais}
\othertitle{Traditional storytellers of the Jequitinhonha Valley (MG): guardians of memory and ancestry: interfaces of oral literature, linguistic variation, and digital technologies}
% if there is a third language title, add here:
%\othertitle{Artikelvorlage zur Einreichung beim Texto Livre Journal}

\author[1]{Cleomar Poletto~\orcid{0009-0008-5689-1507}\thanks{Email: \href{mailto:poletto@reitoria.ufmg.br}{poletto@reitoria.ufmg.br}}}
\affil[1]{Universidade Federal de Minas Gerais, Belo Horizonte, MG, Brasil.}

\addbibresource{article.bib}
% use biber instead of bibtex
% $ biber article

% used to create dummy text for the template file
\definecolor{dark-gray}{gray}{0.35} % color used to display dummy texts
\usepackage{lipsum}
\SetLipsumParListSurrounders{\colorlet{oldcolor}{.}\color{dark-gray}}{\color{oldcolor}}

% used here only to provide the XeLaTeX and BibTeX logos
\usepackage{hologo}

% if you use multirows in a table, include the multirow package
\usepackage{multirow}

% provides sidewaysfigure environment
\usepackage{rotating}

% CUSTOM EPIGRAPH - BEGIN 
%%% https://tex.stackexchange.com/questions/193178/specific-epigraph-style
\usepackage{epigraph}
\renewcommand\textflush{flushright}
\makeatletter
\newlength\epitextskip
\pretocmd{\@epitext}{\em}{}{}
\apptocmd{\@epitext}{\em}{}{}
\patchcmd{\epigraph}{\@epitext{#1}\\}{\@epitext{#1}\\[\epitextskip]}{}{}
\makeatother
\setlength\epigraphrule{0pt}
\setlength\epitextskip{0.5ex}
\setlength\epigraphwidth{.7\textwidth}
% CUSTOM EPIGRAPH - END

% to use IPA symbols in unicode add
%\usepackage{fontspec}
%\newfontfamily\ipafont{CMU Serif}
%\newcommand{\ipa}[1]{{\ipafont #1}}
% and in the text you may use the \ipa{...} command passing the symbols in unicode

% LANGUAGE - BEGIN
% ARABIC
% for languages that use special fonts, you must provide the typeface that will be used
% \setotherlanguage{arabic}
% \newfontfamily\arabicfont[Script=Arabic]{Amiri}
% \newfontfamily\arabicfontsf[Script=Arabic]{Amiri}
% \newfontfamily\arabicfonttt[Script=Arabic]{Amiri}
%
% in the article, to add arabic text use: \textlang{arabic}{ ... }
%
% RUSSIAN
% for russian text we also need to define fonts with support for Cyrillic script
% \usepackage{fontspec}
% \setotherlanguage{russian}
% \newfontfamily\cyrillicfont{Times New Roman}
% \newfontfamily\cyrillicfontsf{Times New Roman}[Script=Cyrillic]
% \newfontfamily\cyrillicfonttt{Times New Roman}[Script=Cyrillic]
%
% in the text use \begin{russian} ... \end{russian}
% LANGUAGE - END

% EMOJIS - BEGIN
% to use emoticons in your manuscript
% https://stackoverflow.com/questions/190145/how-to-insert-emoticons-in-latex/57076064
% using font Symbola, which has full support
% the font may be downloaded at:
% https://dn-works.com/ufas/
% add to preamble:
% \newfontfamily\Symbola{Symbola}
% in the text use:
% {\Symbola }
% EMOJIS - END

% LABEL REFERENCE TO DESCRIPTIVE LIST - BEGIN
% reference itens in a descriptive list using their labels instead of numbers
% insert the code below in the preambule:
%\makeatletter
%\let\orgdescriptionlabel\descriptionlabel
%\renewcommand*{\descriptionlabel}[1]{%
%  \let\orglabel\label
%  \let\label\@gobble
%  \phantomsection
%  \edef\@currentlabel{#1\unskip}%
%  \let\label\orglabel
%  \orgdescriptionlabel{#1}%
%}
%\makeatother
%
% in your document, use as illustraded here:
%\begin{description}
%  \item[first\label{itm1}] this is only an example;
%  % ...  add more items
%\end{description}
% LABEL REFERENCE TO DESCRIPTIVE LIST - END


% add line numbers for submission
%\usepackage{lineno}
%\linenumbers

\begin{document}
\maketitle

\begin{polyabstract}
\begin{abstract}
O presente artigo defende a perspectiva de que um \textit{corpus} de duzentos e cinquenta e cinco narrativas orais, que podem ser ouvidas e lidas no recurso educacional digital “Narrativas em Rede”, criado por este autor, pode ser tipificado como literatura oral, a partir das ideias de \textcite{Cascudo1984}, e os contadores-narradores, concebidos como narradores tradicionais, em consonância com as ideias de \textcite{Benjamin1985}, no ensaio “O narrador: considerações sobre a obra de Nikolai Leskov”. Defende-se a ideia de que os contadores-narradores identificados no Vale do Jequitinhonha entre 1987 e 1997, ao performarem as narrativas orais analisadas, fazem uso de uma oralidade \textit{mista} ou \textit{primária}, tal qual concebida por \textcite{Zumthor1993}, e, ao mesmo tempo, prospectam o saber e a experiência de narradores tradicionais, figura central do ensaio de \textcite{Benjamin1985}. Acredita-se que a criação do recurso digital (site) “Narrativas em Rede” dá voz à oralidade e, ao mesmo tempo, projeta esse \textit{corpus} de literatura oral do Jequitinhonha para dimensões espaciais além do território em que foi gerado. Atualmente, hospedada na Rede Mundial de Computadores, tal literatura torna-se universal, podendo ser acessada de qualquer parte do Brasil e do mundo, por meio de tecnologias digitais como smartphones, tablets, entre outros, consolidando-se como uma ponte entre passado e futuro, unindo tradição e inovação, modos de narrar e modos de socializar das novas gerações.

\keywords{Literatura oral \sep Narradores tradicionais \sep Variação linguística}
\end{abstract}

\begin{english}
\begin{abstract}
This article argues that a corpus of two hundred and fifty-five oral narratives, which can be listened to and read in the digital educational resource “Narrativas em Rede”, created by this author, may be classified as oral literature, based on the ideas of \textcite{Cascudo1984}, and the storytellers, conceived as traditional narrators, in line with \posscite{Benjamin1985} ideas in the essay “The Storyteller”. It is argued that the storytellers identified in the Jequitinhonha Valley between 1987 and 1997, when performing the oral narratives analyzed, make use of a mixed or primary orality, as conceived by \textcite{Zumthor1993}, while, at the same time, drawing upon the knowledge and experience of traditional narrators, a central figure in \posscite{Benjamin1985} essay. It is believed that the creation of the digital resource (website) “Narrativas em Rede” gives voice to orality and, at the same time, projects this corpus of oral literature from the Jequitinhonha region into spatial dimensions beyond the territory in which it was produced. Currently hosted on the World Wide Web, this literature becomes universal, accessible from any part of Brazil and of the world through digital technologies such as smartphones, tablets, among others, consolidating itself as a bridge between past and future, uniting tradition and innovation, modes of storytelling and forms of social interaction of new generations.

\keywords{Oral literature \sep Traditional storytellers \sep Language variation}
\end{abstract}
\end{english}
% if there is another abstract, insert it here using the same scheme
\end{polyabstract}

\section{Introdução}\label{sec-intro}
O artigo faz parte de um recorte de minhas pesquisas desenvolvidas no mestrado na Faculdade de Educação da Universidade Federal de Minas Gerais (FaE/UFMG), cuja temática investigou variações linguísticas em narrativas orais no Vale do Jequitinhonha, em Minas Gerais. Tal pesquisa, já concluída, amplia-se ainda mais no atual contexto, tendo em vista possíveis desdobramentos que estão em investigação no doutorado do Programa de Pós-Graduação em Estudos de Linguagem do Centro Federal de Educação Tecnológica de Minas Gerais (CEFET-MG).

Procura-se defender que um corpus de narrativas orais gravadas no Vale do Jequitinhonha, entre os anos de 1987 e 1997, pode ser concebido como literatura oral, a partir da perspectiva de \textcite{Cascudo1984}, e os contadores destas, concebidos como narradores tradicionais, em consonância com as ideias imanadas do ensaio “O narrador: considerações sobre a obra de Nikolai Leskov”, de \textcite{Benjamin1985}, as ideias de \textcite{Zumthor1993,Busatto2013,Pereira1996} etc.
    
Defende-se que a maioria dos narradores em questão faz uso de oralidade \textit{primária e imediata} como mecanismo de transmissão de conhecimento, sendo eventualmente mediada por saberes decorrentes da escrita e da leitura, tal como preconizado por \textcite{Zumthor1993}.
    
O artigo divide-se em três sessões. Na primeira parte do texto, apresenta-se a concepção de narrador a partir das ideias de Walter Benjamin e estabelecem-se intercessões com as ideias de \textcite{Zumthor1993} sobre linguagem \textit{primária e mista}. Na segunda, disserta-se como dois projetos inventariaram o conjunto das narrativas e narradores no Vale do Jequitinhonha, (MG), entre os anos de 1987 e 1997, utilizando tecnologia analógica de fitas cassete. Apresentam-se, ainda, nessa sessão, alguns fenômenos linguísticos e discute-se como esse \textit{corpus} se constitui em um rico acervo de literatura oral. Por fim, apresenta-se o site “Narrativas em Rede”, que contém um total de 255 narrativas orais que foram convertidas do analógico (fita cassete) para o modelo digital MP3. Atualmente, todo o acervo encontra-se em formato digital (transcrições em PDF) e áudios no formato MP3, podendo ser acessado de qualquer parte do Brasil e do mundo, mediante tecnologias digitais como computador, smartphones, tablets e outros. A conversão do analógico para o digital foi necessária, tendo em vista que o formato analógico já é, há anos, tecnologia obsoleta e se encontra em desuso. A conversão do acervo para novas tecnologias digitais possibilita o acesso a um importante inventário do dialetal em uso no Vale do Jequitinhonha (MG), no contexto de 1987 a 1997. Possibilita, ainda, que escolas e sistemas de ensino se apropriem do dialeto para estudos de variedades linguísticas distintas da variedade de prestígio, veiculada pelos sistemas formais de ensino.

Todas as submissões que contiverem relatos de pesquisas feitas com seres humanos deverão apresentar análise do COEP (apresentação do número CAAE) ou de outros órgãos de ética em pequisa.

\section{O narrador: oralidade e escrita}

Inicia-se a discussão sobre a temática a partir do ensaio “O narrador: considerações sobre a obra de Nikolai Leskov”, de Walter Benjamin. Nesta obra, o autor debruça-se sobre a arte ancestral de contar histórias. A imagem do narrador, figura central no ensaio, representa uma forma de transmissão de conhecimentos, experiências e saberes tradicionais, em contraste com outras formas de transmissão presentes na modernidade. Para \textcite{Benjamin1985}, o narrador é o guardião da experiência e da história humana, detentor de uma sabedoria que é moldada, tecida e transmitida pela experiência. Apropriando-se da linguagem e da oralidade, ele é capaz de tecer ensinamentos na teia de seus relatos. Todavia, para ele, tal arte encontra-se em declínio diante das profundas transformações sociais e do advento de outras formas de transmissão.

Segundo o autor, o “conselho” (referindo-se às mensagens transmitidas pelas narrativas orais) “tecido na substância viva da existência tem um nome: sabedoria” \cite[p. 200]{Benjamin1985}. O que é falado é tecido, é construído, é vivido. Nessa tecitura, a palavra do narrador reporta fatos vividos. Fala porque viu, fala porque viveu.
    
Destaca, ainda, que “a arte de narrar está definhando porque a sabedoria – o lado épico da verdade – está em extinção”. Esse processo “vem de muito longe […] nada seria mais tolo que ver nele um sintoma de decadência ou uma característica moderna” \cite[p. 201]{Benjamin1985}.

Para o autor, o processo que expulsa gradualmente a narrativa da esfera do discurso vivo, aquele ancorado na experiência, vem ocorrendo ao longo do tempo, com a evolução secular de forças produtivas e de inventos industriais. O surgimento do romance e a invenção da imprensa são os primeiros indícios da morte do narrador e das narrativas, já no início do Período Moderno.

No ensaio “O narrador: considerações sobre a obra de Nikolai Leskov”, emergem dois tipos arcaicos de narradores que, metaforicamente, prospectam os demais: o narrador-viajante – “quem viaja tem muito a contar” –, este associado à figura do marinheiro comerciante, e o narrador “camponês sedentário” \cite[p. 198]{Benjamin1985}, que vive da terra, que tem raízes e histórias no seu território.

\textcite[p. 17-18]{Busatto2013} considera que o narrador tradicional – este do Vale do Jequitinhonha (MG) – “narra para se sentir vivo […]. Esse personagem e suas palavras aladas sempre estiveram presentes na alma da comunidade”. Ainda, de acordo com a autora:

\begin{quote}
    Era um sujeito que se valia da narração oral como via para organizar o caos, perpetuar e propagar os mitos fundacionais das suas culturas. Um sujeito que mantinha vivo o pensamento do seu povo por meio da memória prodigiosa e que o divulga por meio da arte \cite[p. 18]{Busatto2013}.
\end{quote}

Em \textit{A arte de contar histórias no século XXI}, Busatto relata que a figura do narrador-contador vem sofrendo alterações no decorrer da história, não em sua opção pela palavra ou pela oralidade. Pondera que esses herdeiros da tradição oral “já se encontram inseridos num contexto mediado pelos novos meios de comunicação e transmissão de saber” \cite[p. 19]{Busatto2013}.

\textcite[p. 18]{Zumthor1993}, em \textit{A letra e a voz}, destaca três “tipos” de oralidades presentes nas sociedades, que correspondem a “três situações de cultura”. A primeira é a oralidade \textit{primária e imediata}, na qual os narradores não tem nenhum contato com a escrita. Para o autor, esse tipo de oralidade “se encontra apenas nas sociedades desprovidas de todo sistema de simbolização gráfica, ou nos grupos sociais isolados e analfabetos”.
    
A segunda é a oralidade \textit{mista}, em que os sujeitos já convivem com a escrita, embora esta pouco influencie em seu cotidiano. “Denominei […] oralidade mista, quando a influência do escrito permanece externa, parcial e atrasada […] procede da existência de uma cultura escrita, no sentido de possuidora de uma escritura”. Por fim, há a oralidade \textit{segunda}, que é característica da “cultura letrada, na qual a expressão é marcada mais ou menos pela presença da escrita” \cite[p. 18]{Zumthor1993}.
     
Defende-se que os contadores-narradores do Vale do Jequitinhonha, ao performarem as narrativas orais, fazem uso de uma oralidade \textit{mista ou primária}, tal qual concebida por esse autor, ao mesmo tempo em que prospectam o saber e a experiência de narradores tradicionais, figura central do ensaio de \textcite{Benjamin1985}, em consonância com as discussões de \textcite{Busatto2013,Pereira1996}, entre outros.
    
É sobre essas narrativas e esses narradores, prospectados pelo filósofo Walter Benjamin, que nos propomos a falar e como suas vozes podem ser ouvidas, mesmo sendo ecos de tempos passados, de tempos de mudanças, muitas já ocorridas, outras ainda por ocorrer.

\section{Narradores(as) e narrativas: uma busca}
Em 1987, inicia-se, no Vale do Jequitinhonha, (MG)\footnote{O Vale do Jequitinhonha (MG) é uma mesorregião do estado de Minas Gerais, Brasil, com uma área de aproximadamente 65 mil km². É composto por 52 municípios e é cortado pelo Rio Jequitinhonha, que dá nome a região. A formação histórica do território do Vale remonta às primeiras entradas em busca de ouro e pedras preciosas à época do Brasil Colônia (cf. \textcite{Poletto2025}).}, a busca por narradores(as) e por narrativas orais. Nesse contexto, os pesquisadores Reinaldo Martiniano Marques, da UFMG, e Vera Lúcia Felício Pereira, da Pontifícia Universidade Católica de Minas Gerais, (PUC Minas), iniciaram o projeto “Literatura Oral do Vale do Jequitinhonha”. Esse projeto contou com financiamento do Conselho Nacional de Desenvolvimento Científico e Tecnológico (CNPq) e da PUC Minas.

A pesquisa, de acordo com \textcite[p.~18]{Pereira1996},“consistiu, na primeira fase, em ganhar a aquiescência dos contadores-narradores para depois conversar com eles e poder ouvi-los”. Segundo a autora, “tomavam-se notas, gravava-se e faziam-se os registros, somente com o consentimento destes e se isso não os desagradasse” \cite[p. 18]{Pereira1996}. Foi assim, entres idas e vindas, que esses pesquisadores foram conhecendo, ganhando confiança e abrindo caminho para novos retornos e registros futuros.

 \textcite{MarquesPereira1988}, em artigo publicado no periódico “O Eixo e a Roda”, detalham os objetivos e metas do projeto: “trata-se de uma pesquisa sobre a literatura popular do Vale, empenhada na recolha e análise de contos, casos, histórias e lendas da região” \cite[p. 173]{MarquesPereira1988}. Envolve, segundo os pesquisadores, o problema da memória, tendo por justificativa a intensa e relevante produção cultural da cultura popular do Vale do Jequitinhonha, a preservação da identidade cultural da região e a criação de uma relação mais efetiva entre sociedade e universidade.
 
Ainda de acordo com os autores, o projeto objetivava identificar contos no Jequitinhonha; caracterizar seus processos de produção, circulação e recepção; determinar as recorrências temáticas e as preocupações formais neles presentes; e divulgar a literatura popular do Vale, partindo inicialmente de cinco municípios – Diamantina, Serro, Minas Novas, Turmalina e Araçuaí – posteriormente expandindo-se para outras cidades.

Ao término das atividades de campo, em 1990, o projeto registrou, em cinco municípios, o total de 177 narrativas que podem ser tipificadas como literatura oral\footnote{Segundo \textcite{Cascudo1984}, a denominação “literatura oral” é de 1881 e foi cunhada por Paul Sébillot: “Sua característica é a persistência pela oralidade. A fé é pelo ouvir, ensina São Pedro” \cite[p. 23]{Cascudo1984}. Quanto às fontes, há duas que mantêm viva a corrente da literatura oral. “Uma é exclusivamente oral, resume-se na estória, no canto popular e tradicional, nas danças de roda, danças cantadas, danças de divertimento coletivo […] contos, anedotas, adivinhações, lendas, etc” \cite[p. 23]{Cascudo1984}. A outra, segundo o autor, se dá pela reimpressão dos antigos livrinhos vindos da Espanha ou de Portugal, que logo são absorvidos pela oralidade. Esse processo pode ser revisitado a partir das ideias de \textit{oralidade mista} descrita por \textcite{Zumthor1993}, em que pessoas convivem com a escrita, porém o sistema pouco influencia o seu cotidiano. Os narradores do Vale do Jequitinhonha faziam uso da leitura, ao menos alguns. Todavia, quase a totalidade eram analfabetos ou com baixa escolaridade, não tendo a leitura como mecanismo para transmissão do conhecimento
(cf. \textcite[p. 67-80]{Poletto2025}).} – contos de encantamento, ensinamento, magia, “causos” e lendas do folclore brasileiro e da região. Nessa etapa, foram identificados 42 contadores(as), sendo 16 mulheres e 27 homens.

A partir das gravações realizadas pelo projeto .“Literatura Oral”, surgiu a obra \textit{O artesão da memória do Vale do Jequitinhonha}, de Vera Lúcia Felício Pereira, publicada em 1996 pela Editora UFMG e pela Editora PUC Minas. Nesse livro, a autora retoma os argumentos de Walter Benjamin sobre os narradores(as) e sobre suas narrativas, “causos” e contos:

\begin{quote}
    O conto maravilhoso, há séculos transmitido oralmente, exerce uma função social. Seu enredo oferece conselhos, auxílios, e preenchendo carências satisfaz expectativas dos que leem e dos que não leem… e deixa entrever aos ouvintes e leitores formas mais satisfatórias de vida e o modo de alcançá-las \cite[p. 51]{Pereira1996}.
\end{quote}

Por meio dessa obra, a autora apresenta um pouco da vida de alguns narradores-contadores e registra, a partir da narrativa oral, um pouco da história do próprio Vale do Jequitinhonha. Neste sentido, a palavra do narrador constitui-se em fonte de informação e registro de tempos passados, configurando-se como memória e história de um povo.

Segundo \textcite[p.~20]{Pereira1996}, o narrador do Vale “é tropeiro, como João de Deus, 87 anos quando gravado, quarenta e dois vividos como tropeiro”. As narrativas de João de Deus Vaz relatam histórias de tropeiros, de massas de tropas e da logística de transporte disponível nos interiores de Minas e do Brasil à época. De acordo com João de Deus:

\begin{quote}
    A tropa saía daqui, a tropa num tinha massa grande. Só tropa pequena de três légua. A massa maior de tropa era quatro légua. Daí pra trás. Daqui a Teófilo Otoni são doze ponto de arrancha. De volta a mesma coisa. Doze \footnote{Em todas as inserções de fragmentos de narrativas orais dos(as) narradores(as), foi mantido o texto conforme a transcrição, com todas as variedades linguísticas.} dia pra cá, doze dia pra lá \cite[p. 21]{Pereira1996}.
\end{quote}

João de Deus Vaz foi gravado em 1987, em Minas Novas. Nascido em 1901, certamente viveu à época em que o principal meio de transporte deste país era a massa de tropas. Viveu a transição de dois modais, se assim podemos dizer: o modal puxado por tração animal (mulas e burros) e o modal mecânico rodoviário. Ele testemunhou, portanto, o fim de um meio de transporte que levou progresso e desenvolvimento para o interior do país. Sr. João narra a transição dos modais na narrativa “A história do carro”, gravada em cassete, em 1987, por Reinaldo Martiniano Marques. O relato conta a chegada do primeiro carro no Vale. É um testemunho de quem viveu para ver e viu como a estrada, antes utilizada por tropas, foi sendo lentamente adaptada para a circulação de veículos:

\begin{quote}
    E de lá, de Sete Porta, de Concórdia até a Situbinha é um trecho de istrada muito ruim. Aí eles entrava po córrego do Tamburi, entrô po córrego do Tamburi, passano pela fazenda do Heleno, pra saí na Muvuca, pra saí na Muvuca. Aí uns trecho que era lugar que tinha quatro home de chave, picareta e pá. Aqueles lugá, pa podê passá, fazia aquelas tauba pra Manoel i podeno dirigi. Otros lugá eles vinha fazeno tábua, naqueles lugá de brejo, brejado, mole, eles forrava co’ aquelas tábua, o cavalo vinha, era um carro pequeno, num era grande não, o carro passava (Vaz, 1987)\footnote{O fragmento citado acima faz parte da narrativa oral gravada em 1987 em Minas Novas (MG)  e pode ser ouvida, pois encontra-se no formato MP3 no site \url{https://drive.google.com/drive/folders/1N6P0rL_0bcnVtDb4R11GAqnZKPANsIjZde} de recurso educacional “Narrativas em Rede”; link para o áudio: A transcrição também se encontra disponível e pode ser lida simultaneamente à escuta do áudio: \url{https://drive.google.com/drive/folders/1D55VLM3eMXLDyRAD24CG7F51glUv7NyU.}}.
\end{quote}

O contador é agricultor, lavrador, boia-fria, como Neném de Araçuaí, um itinerante que retorna com outros saberes. Por analogia, Neném simboliza o “marinheiro mercador” descrito por \textcite{Benjamin1985}. É o narrador que sai de sua terra, no contexto do Vale, como migrante e boia-fria para trabalhar em lavouras de cana-de-açúcar, café e laranja, nos estados de Minas Gerais e São Paulo, configurando o fenômeno do êxodo rural muito comum no Vale. Neném e outras milhares de pessoas do Vale simbolizam o que \textcite{Silva1999} denomina de massa de errantes em busca de trabalho e sobrevivência. São pessoas que perderam, na região, as condições de sobrevivência decorrentes da expansão das fronteiras agrícolas, política implementada pela Ditadura Militar brasileira, na década de 60/70.

Vicente Nica, neste contexto, é outro narrador que simboliza a perda da terra e das dimensões simbólicas que ela possuía. Mas, resiliente, pôs-se em luta. Desse modo, Vicente Gonçalves Afonso – mais conhecido como Vicente Nica, nascido em 1932, na Grota do Barreiro, em Turmalina (MG), posseiro de terras posteriormente reivindicada por fazendeiros – não se tornou Vicente Nica apenas pela forma como lavrava a terra enquanto agricultor, mas enquanto agricultor que lutou e denunciou a expropriação das terras no ciclo desenvolvimentista empreendido na década de 60/70.

É histórica a luta travada por Vicente Nica e outros posseiros de Mato Grande e São Miguel para manter suas posses quando a família França reivindicava para si a propriedade de 14 mil hectares de terra, já ocupados por lavradores que, de geração em geração, faziam uso delas \cite{Leite2015}. Além disso, Vicente Nica denunciou a ocupação das chapadas pela monocultura do eucalipto, que se alastrava pelo Vale, encurralando os agricultores nas grotas, em virtude da expropriação das chapadas, que eram áreas de uso comum. Destacou-se, ainda, por conceber um modelo de agricultura que buscava consorciar o uso do solo com a proteção do meio ambiente.

Na obra de \textcite{Silva1999}, encontra-se a denúncia de Vicente Nica à forma de ocupação das terras pela monocultura de eucalipto, que no Vale se deu com o apoio técnico e financeiro do Estado:

\begin{quote}
    Eles compravam as terras do lado, iam cercando o sujeito. Plantava eucalipto e cercavam a área. Eles queriam terra plana. Passava com o trator nas terras do sujeito que não queria vender, e ai iam destruindo tudo. Foi o jeito de forçar a venda. Vendia pelo preço que eles queriam […] ninguém entendia de valor de terra \cite[p. 59]{Silva1999}.
\end{quote}


O narrador é garimpeiro, agricultor, poeta, escritor; é fonte de inspiração. É trabalhador no serviço público dos correios como Pedro Braga, que, para não deixar cair no esquecimento as lembranças do Vau, passou a registrar tudo em papel. \textcite{Pereira1996} descreve:

\begin{quote}
No seu objetivo de não deixar que sua terra caia no esquecimento total, ele registra cuidadosamente nos seus cadernos toda a história do lugar. As lendas referentes aos pontos topográficos, as famílias proprietárias e sua descendência e os casos dos escravos […] os cadernos de Seu Pedro são preciosidades à espera de estudiosos que, admitindo o espaço da diversidade e do conflito, sempre presentes, queiram dar-lhes agora a voz \cite[p. 21]{Pereira1996}.
\end{quote}

\textcite{Souza2006} foi um desses estudiosos mencionados pela autora. Com seus trabalhos, contribuiu para dar voz a Pedro Braga. Localizou, junto à família do narrador, quatorze cadernos manuscritos, em verso e prosa, muitos dos quais integram suas pesquisas de mestrado e doutorado. Ainda de acordo com \textcite{Souza2018}, Pedro Braga inspirou Eliane Caffé a criar o personagem Antônio Biá, do filme \textit{Narradores de Javé}, lançado em 2003. A cineasta tomou conhecimento de Pedro Braga por meio do livro \textit{O artesão da memória no Vale do Jequitinhonha}, de \textcite{Pereira1996}, e com ele trocou correspondência\footnote{Eliane Caffé tomou ciência da existência de Pedro Braga enquanto gravava o filme Kenoma \cite{kenoma1998}, em Araçuaí (MG), na comunidade do Itira, que está localizada às margens do Rio Jequitinhonha e do Rio Araçuaí, onde as águas se encontram. Para mais informações, acesse: \url{https://www1.folha.uol.com.br/fsp/ilustrad/fq2912199917.htm}}.

Sobre Pedro Braga, \textcite{Pereira1996} escreve:

\begin{quote}
    Seu Pedro é, no Vau, agente dos Correios e, ao sentir-se na obrigação de manter o posto aberto, único fio de comunicação e vínculo com o progresso, escreve cartas aos que conhece ou de quem já ouviu falar. Como muitas pessoas o procuram, atraídas pela sua fama de contador, solicita aos visitantes o endereço na esperança de que, enviando e recebendo muitas cartas, consiga impedir o fechamento da única instituição que liga sua terra ao resto do país \cite[p. 30]{Pereira1996}.
\end{quote}

Se Neném de Araçuaí prospecta a imagem do narrador “marinheiro mercante”, Pedro Braga prospecta a figura do “narrador camponês”, sedentário, que, preso à sua terra, mantém, por meio da escrita, a memória do passado e suas tradições. Os(as)  narradores(as) do Vale, pensados no par viajante/sedentário, marinheiro/camponês, são figuras que não estão em oposição; ao contrário, complementam-se num devir. Formam e se conformam.

\textcite{Benjamin1985} desenvolve tal pensamento, no ensaio “O Narrador: considerações sobre a obra de Nikolai Leskov”, dizendo que os camponeses e os marujos foram os primeiros mestres da arte de narrar; no entanto, foi da síntese entre o sedentário e o viajante que a arte de narrar foi aperfeiçoada:

\begin{quote}
    Se os camponeses e os marujos foram os primeiros mestres da arte de narrar, foram os artífices que a aperfeiçoaram. No sistema corporativo associava-se o saber das terras distantes, trazidos para casa pelos migrantes, com o saber do passado, recolhido pelo trabalhador sedentário \cite[p. 199]{Benjamin1985}.
\end{quote}

Realizada a apresentação de alguns narradores, passa-se agora para a segunda etapa, discutindo e apresentando o trabalho realizado pelo segundo projeto que se pôs a registrar narradores(as) e narrativas no Vale do Jequitinhonha (MG).

\subsection{Narradores(as) e narrativas: novos encontros}

O Boletim da UFMG nº 1321, do ano de 2001, relata que o projeto “Quem Conta Um Conto Aumenta Um Ponto” foi criado no ano de 1995 pela professora Sônia Queiroz, da Faculdade de Letras da UFMG. Durante sua vigência, foi financiado pela Pró-Reitoria de Extensão da UFMG (PROEX). Destaca-se, todavia, que a execução se deu dentro do programa “Polo Jequitinhonha da UFMG”\footnote{Existente desde 1996, o Polo Jequitinhonha da UFMG desenvolveu e desenvolve inúmeros projetos em diferentes áreas: educação, meio ambiente, geração de renda, direitos humanos, direitos das mulheres, preservação e gestão das águas, valorização e desenvolvimento da cultura regional, artesanato, música, canto, dança e literatura oral etc.}.

Sob o guarda-chuva do “Polo Jequitinhonha”, esse projeto recebeu da professora Vera Lúcia Felício Pereira, por meio de doação\footnote{As narrativas orais, compostas por áudios e suas respectivas transcrições, foram convertidas para o formato MP3 e encontram-se disponibilizadas no site “Narrativas em Rede”, criado por este autor, com o objetivo de preservar, organizar e difundir esse patrimônio linguístico-cultural. Já o acervo físico, constituído por fitas cassete que somam centenas de horas de gravação, encontra-se depositado na Biblioteca Universitária da UFMG, integrando a coleção “Marizinha Nogueira/Polo Jequitinhonha da UFMG”. Tal acervo representa importante fonte documental para pesquisas nas áreas da sociolinguística, memória oral, cultura popular e estudos da linguagem no contexto do Vale do Jequitinhonha.}, o acervo de narrativas orais gravadas pelo projeto “Literatura Oral” entre 1987 e 1990. De posse do inventário com 177 narrativas e 42 narradores(as) identificados(as), o projeto “Quem Conta Um Conto Aumenta Um Ponto” retoma as gravações em 1995 e, em 1997, termina seus registros em campo. Assim, em mais de dois anos de pesquisa, ele inventariou e gravou mais 78 narrativas e 13 contadores, 7 homens e 6 mulheres, nos municípios de Araçuaí, Minas Novas e Turmalina. A pesquisa expandiu-se para Jenipapo de Minas e Malacacheta, municípios que não haviam sido alcançados pelo projeto anterior, de 1987 a 1990. Observa-se que um contador do município de Minas Novas, Sr. Joaquim Soares Ramos, foi gravado pelos dois projetos.

Desta forma, ao longo de 10 anos de pesquisas, os dois projetos conseguiram registrar, documentar e gravar um total de 255 narrativas e 55 narradores(as). Tais narrativas podem ser tipificadas como um acervo de literatura oral, contos de encantamento, ensinamento, magia, “causos” e lendas do folclore brasileiro e da região.

Se ao projeto “Literatura Oral” coube gravar 69,4\% das narrativas, ao projeto “Quem Conta Um Conto Aumenta Um Ponto” coube ampliar a busca por novos contadores(as) e novas narrativas, e, ao mesmo tempo, transcrever\footnote{A transcrição pode ser conceituada como técnica de escrita da fala. Tudo que é falado é transformado em palavras fielmente, sem acréscimos, recortes ou supressões. Procura-se aproximar ao máximo a escrita da fala. Por exemplo, para expressar dúvidas, silêncios e hesitações em um processo de transcrição, usa-se reticências […]; risos devem ser identificados com a palavra “risos” entre parênteses.} duzentas 200 narrativas dentre as 255 gravadas. Alguns trabalhos foram publicados pelo “Polo Jequitinhonha UFMG”, a partir das transcrições; a saber: \textit{7 histórias de encanto e magia} (livro e CD), de \textcite{Queiroz1999}; \textit{No tempo em que os bichos falavam} (livro e CD), de \textcite{Queiroz2009}; além de outras obras que, direta ou indiretamente, tiveram como substrato de pesquisa as narrativas orais gravadas e transcritas, como o \textit{Dicionário do dialeto rural no Vale do Jequitinhonha}, de \textcite{AntunesFerraz2013}, e o audiolivro \textit{Negros pelo Vale}, de \textcite{Souza2014}.

\textcite{Souza2014} analisa 35 narrativas e esclarece que a pessoa negra é representada, na maioria das vezes, como escravizada, tendo em vista que foi essa a condição que marcou sua entrada no Vale do Jequitinhonha para exploração do ouro e pedras preciosas, em Diamantina, Minas Novas e demais municípios da região. Em outras narrativas, embora não se figurem como escravizadas, elas são situadas em uma posição inferior na hierarquia social, sendo representadas, em alguns casos, sob um viés estereotipado.

\textcite{AntunesFerraz2013}, com colaboradores, registraram, a partir das narrativas, mais de 1.000 unidades léxicas do dialetal da região do Vale. A obra registra a lexicografia da cultura oral no Vale e traz conhecimentos linguísticos, históricos e sociais. Segundo Maria Cândida Trindade Costa Seabra, prefaciadora do dicionário, a produção da obra “dá existência escritural à fala” \cite[p. 11]{AntunesFerraz2013}, evidenciando o movimento de transposição da oralidade para o registro escrito e seu valor para preservação e análise da língua em uso.

A partir das pesquisas realizadas no Mestrado em Educação e Docência da Faculdade de Educação da UFMG, foi possível constatar que as 255 narrativas em áudio e as 200 narrativas transcritas formam um \textit{corpus} de mais de 168 mil palavras. Esse conjunto de áudios (ainda em fita cassete até 2024) e transcrições não havia passado por uma investigação sob a perspectiva da sociolinguística variacionista (pelo menos, não totalmente). Ressalva-se ao argumento explicitado o trabalho de transcrição realizado pelo projeto “Quem Conta Um Conto Aumenta Um Ponto”, que, de forma criteriosa, conseguiu materializar na escrita a variedade linguística falada na década de 1987 a 1997, preservando as marcas da oralidade dos narradores. Assim, a escritura da língua em uso, registrada pelas transcrições, com suas singularidades e particularidades, constituiu-se em um acervo de grande valor para pesquisas na perspectiva da sociolinguística variacionista, tanto para estudos sincrônicos quanto para estudos diacrônicos.

Esse conjunto de áudios e transcrições, tipificado genericamente como contos, encontra-se disponível no recurso educacional digital “Narrativas em Rede” (disponível em: \url{https://www.ufmg.br/polojequitinhonha/narrativas-em-rede/).} As narrativas podem ser utilizadas por meio de tecnologias digitais móveis por professores, pesquisadores e alunos, tanto para usos educativos, tanto para usos estéticos, visto que possuem a “isca” do encanto e da magia, típica do discurso narrativo.

\subsection{Examinando alguns fenômenos linguísticos do Vale do Jequitinhonha}

A investigação de uma variedade linguística, no âmbito da sociolinguística, dá-se pela descrição detalhada de suas singularidades. Em geral, parte-se de uma análise descritiva do objeto de estudo. O termo técnico para pôr em evidência um fenômeno linguístico variacionista denomina-se \textit{envelope de variação}. Esse conceito é empregado nas pesquisas para descrever, de forma minuciosa, uma variável, as variantes e os contextos em que podem ou não ser produtivas. \textcite{coelho2012sociolinguistica} pontuam que, para a descrição do \textit{envelope de variação}, faz-se necessário reconhecer o fenômeno em mudança, descrever a variedade e identificar as variantes. Assim, deve haver equivalência de sentido na existência de duas ou mais variedades.
 
Para \textcite[p.~68]{BortoniRicardo2019}, “as formas que supostamente transmitem o mesmo conteúdo semântico, expresso com recursos linguísticos distintos, vão caracterizar regras variáveis, e suas alternativas são denominadas variantes”. Ou seja, as formas “nós fomo”, “nós fomos”, “nós foi” e “nós fumu” são equivalentes, pois produzem o mesmo sentido.

Apresentam-se, a seguir, algumas variedades extraídas do fragmento da história “O Cueio e o Cascavelo”, narrativa gravada em 1997 pelo projeto “Quem Conta Um Conto Aumenta Um Ponto”, no município de Minas Novas, com o contador Joaquim Soares Ramos. O áudio pode ser ouvido no endereço \url{https://drive.google.com/file/d/1jpoV3oAQZ6PmXvFCWf6rlqxu05abVhde/view}. A leitura da transcrição pode ser simultânea com o áudio no link \url{https://drive.google.com/file/d/1aJKsS9PkgNTUPUU-DjAa_cTvJZbC5fjZ/view.}. 

A narrativa conta a história de um “cueio”, um “cascavelo”, uma onça e “uma” tigre que, ao compartilharem o mesmo espaço, inserem-se em um cenário mágico de conflitos e artimanhas. O cueio e o cascavelo são personagens centrais na narrativa e o ápice do enredo gira em torno deles. Tendo em vista que a história é longa, apresentam-se, a seguir, apenas alguns fragmentos da narrativa. No entanto, recomenda-se a leitura integral da história para melhor compreensão.

\begin{quote}
    [Havia um \textit{cascavelo}...] Então, um dia, a onça, \textit{andano}, \textit{chegô}, deu c’aquele lugar, \textit{aquê} limperão \textit{dibaxo} daquela lapa lá, \textit{ês foi chegano} assim, \textit{entrano} assim \textit{d’ u’a vez}, mas quando ela deu com a vista no \textit{cascavelo}, ela foi e \textit{parô}. Foi, \textit{priguntô ele}, \textit{falô} [...]
    
— Ô \textit{ti cascavelo}, isso aqui é seu? [...]

— Ó, \textit{cê} podia dá eu um \textit{lugazin d’eu ficá} aqui? Tem muita distância […]

[E o cascavelo responde:] 

— Ó, \textit{cê} fica daqui pra lá, ó, \textit{cê} pode \textit{andá} ó, mas desse pedaço aqui pra cá \textit{cê} num anda não que se \textit{ocê passá} daqui pra \textit{dento}, eu \textit{mat’ocê}.

Então a onça ficava, a onça ficava \textit{naquê} pedaço ali e andava, andava... De noite o cascavel também saía pa ‘\textit{rumá} caça. Andava, andava, \textit{vortava}, tá lá \textit{dibaxo} da lapa lá.

Um dia \textit{envém} a tigre. \textit{Chegô, falô}:

— \textit{Ah, sá} onça! Mas que \textit{lugazin} bom, casão que \textit{cê arrumô} aqui, \textit{sá} onça! 

\textit{Falô}:

— Não, isso aqui é de \textit{ti cascavelo}, moço. 

\textit{Falô}:

— Eu \textit{quiria ficá qui} mais \textit{ocê}. […]

\textit{Cascavelo maircô} lá, \textit{falô}:

— Ó, aí ó, daqui pra lá ó, \textit{cês} pode \textit{andá} todo canto, mas desse pedaço pra \textit{dento} aqui \textit{cês} num passa não que se \textit{passá ieu} mato.

Então \textit{ês tá} lá. De noite andava. Andava a onça mais a \textit{tigre}. O \textit{cascavelo tamém} saía, andava, \textit{vortava} pra lá.

Um dia \textit{chegô} o \textit{cuei co aque’as} patinha, ‘\textit{ducado}. \textit{Chegô, olhô lá dibaxo} e \textit{viu ‘quele lugazão lá, chegô}: — Bença, tia onça! Bença, \textit{tia tigre}! Bença, \textit{ti} cascavel! [...]

— Ô, \textit{sá} onça, cê \textit{pudia dexá} eu \textit{ficá} aqui mais \textit{ocê}. [...] \textit{cê arranjô} uma casa boa aqui.

Ora, \textit{falô}:

— Não, pede \textit{ti cascavelo} que aqui é dele, né nosso não.

Aí \textit{falô}, \textit{falô}:

— Ô \textit{ti cascavelo}, \textit{cê pudia dexá} eu \textit{ficá} qui mais \textit{sá} onça, tigue...

\textit{Falô}:

— \textit{Cuei, pó ficá}. Olhe, \textit{cuei}, cê fica aqui assim, ó, nesse \textit{lugazin} aqui, \textit{cuei}. Mas já num \textit{falô co} ele que se ele passasse pra \textit{dento}, que \textit{matav’ele} não \textit{poquê} é bem tratado (grifos nossos).
\end{quote}

Analisando os fragmentos da narrativa, podem-se verificar diversos fenômenos linguísticos em variação, dentre os quais se destacam:

\begin{enumerate}[label=\alph*.]
    \item \textit{Metaplasmos}: são alterações fonéticas que ocorrem em palavras com a evolução da língua. Podem ocorrer por acréscimo, no início ou fim de léxicos, por substituição no interior de palavras, permuta, apagamento etc.
    
    No metaplasmo por aférese, ocorre o apagamento ou a redução de segmento no início de palavras. São exemplos, na narrativa, formas como \textit{ocê, cê} (“você”), \textit{‘ducado} (“educado”), \textit{‘quele} (“aquele”), \textit{qui} (“aqui”), e \textit{‘rumá} (“arrumar”) e \textit{tá} (“está”).
    
    No metaplasmo por apócope, ocorre a subtração de um fonema no final da palavra. Na narrativa, observa-se a perda da consoante /r/ no final de léxicos, o que pode apontar para a perda do infinitivo tanto em verbos isolados quanto em construções com verbo auxiliar juntamente a um principal, como em: “Eu \textit{quiria ficá} qui mais ocê” e “Ô ti cascavelo, cê \textit{pudia dexá} eu ficá qui mais sá onça, tigue…”. Nessas duas construções, identificam-se os verbos auxiliares (“querer” e “poder”) e os verbos (“ficar” e “deixar”) como principais. Ainda no âmbito dos metaplasmos, evidencia-se apócope em formas como \textit{lugazin} (“lugarzinho”), \textit{co} (“com”), \textit{ti} (“tio”), \textit{aquês} (“aqueles”), \textit{ê} (“ele”) e \textit{ês} (“eles”). Ocorrem, ainda, metaplasmos por acréscimo (paragoge) na palavra \textit{cascavelo} (“cascavel”), e metátese, na palavra \textit{priguntô} (“perguntou”).
    \item \textit{Apagamento de ditongos}: para \textcite [p. 101]{MattosoCamaraJr1978}, o ditongo é uma sequência de duas vogais pronunciadas em uma mesma sílaba. Nos fragmentos, destacam-se os apagamentos dos ditongos \textit{/ai/}, \textit{/ei/} e \textit{/ou/}. Para o ditongo /ai/, tem-se como exemplo a palavra \textit{dibaxo} (“debaixo”). Para o ditongo /ei/, aponta-se o verbo \textit{dexá} (“deixar”). No ditongo /ou/, observam-se as formas verbais regulares da primeira conjugação terminadas em “ar”, quando conjugadas na 3ª pessoa do singular do pretérito perfeito do indicativo, como \textit{priguntô, chegô, parô, falô e arrumô}, além do substantivo \textit{oto} (“outro”).
    \item \textit{Rotacismo}: o fenômeno é definido pela troca da líquida lateral \textit{/l/} pela líquida vibrante \textit{/r/}. O rotacismo ocorre, em geral, em sílabas com um encontro de consoantes e vogal, denominadas CCV (consoante-consoante-vogal), como nas palavras “bloco” (por \textit{broco}) e “planta” (por \textit{pranta}). Todavia, pode ocorrer também em outros contextos com palavras com /l/, como “alça”, “calça”, “falta” e “altura”, rotacionando para a vibrante /r/ como nas palavras \textit{arça}, \textit{carça}, \textit{farta}, \textit{artura} e \textit{farsa}. Na narrativa, o fenômeno aparece diversas vezes, como em: “Então ‘rastô, jugô fora, \textit{vortô} o cuei logo inventô fazê um \textit{armoço} lá. [...] Ô, ti cascavel, fazê um \textit{armoço} aí pra nós.”.
    \item \textit{Despalatização}: é definida como a perda da pronúncia no palato alto e a ocorrência da pronúncia no palato médio. Perda do traço palatal na articulação de um fonema. Na variedade de prestígio, a palavra “coelho” se pronuncia no palato alto. Com a perda do traço palatal \textit{/lh/}, a pronúncia se dá no palato médio. Na narrativa, observa-se a ocorrência de \textit{cuei} (11 vezes), \textit{cueio} (7 vezes) e \textit{cuelho} (2 vezes), ou seja, três variantes para designar o animal coelho.
    \item \textit{Degeminação do gerúndio}: trata-se do apagamento da oclusiva dental /d/ no morfema \textit{/ndo/}, que marca o gerúndio nos verbos. Na variedade de prestígio, o gerúndio é realizado pelo encontro consonantal \textit{/nd/} em verbos. Na variedade analisada, constata-se o apagamento dessa consoante, como no fragmento: “O cueio fez, arrumô um tacho de cobre, ‘ranjô uma trempe, colocô o tacho em riba, foi \textit{carregano} água n’uma vasia pequena. Foi \textit{pono} no tacho até que ele encheu o tacho d’água. [...] ê pôs fogo. [...] Pôs fogo, tá só \textit{tiçano} fogo. A panela tá \textit{queimano}, ê tá \textit{tiçano} fogo”.
    \item \textit{Alçamento de vogal}: o alçamento é um processo fonológico que consiste em elevar as vogais médias \textit{/o/} e \textit{/e/}, que passam a ser realizadas com pronúncia alta, com /u/ e /i/, respectivamente. De acordo com \textcite[p. 48]{BortoniRicardo2011} , “a literatura filológica tradicional no Brasil tem considerado a elevação das vogais pretônicas como consequência da regra de harmonia vocálica, isto é, a vogal pretônica é assimilada à vogal tônica alta seguinte”. O fragmento apresenta os seguintes exemplos: “Aí o cascavelo falô: — mas que \textit{minino} bem \textit{insinado} esse \textit{minino}” e “Um dia, quando a onça... às vez a tigue ia saí primero, esbarrava na onça, a onça dav’um rosnado e o cuie tá pono \textit{sintido} naquilo. \textit{Sintido}.”.
    \item \textit{Alteração pronominal}: na narrativa, a alteração de alguns pronomes aparece com frequência, como no caso dos pronomes do caso reto “ele(s)”, os pronomes demonstrativos “aquele(s)” e do pronome de tratamento “você”, conforme se observa nos fragmentos: “Então, um dia, a onça, andano, chegô, deu c’aquele lugar, \textit{aquê} limperão dibaxo daquela lapa lá, \textit{ês} foi chegano assim [...]”; “— Ó, \textit{cê} pudia dá eu um lugazin d’\textit{eu} ficá aqui? [...]”; “— Ó, \textit{cê} fica daqui pra lá, ó, \textit{cê} pode andá ó, mas desse pedaço aqui pra cá \textit{cê} num anda não que se ocê passá daqui pra dento, eu \textit{mat’ocê}. […]”; “Quando encheu o tacho d’água, ê pôs fogo. [...]”; “[...] \textit{ê} tá tiçano fogo” e “começô a frevê, \textit{ê} já […]”.
\end{enumerate}
    
A análise acima demonstra variações linguísticas encontradas no Vale do Jequitinhonha (MG). Entre os achados, destacam-se fenômenos como o apagamento do ditongo \textit{/ou/}, com índices de 95,6\% em verbos, ante 4,40\% na variedade de prestígio. Em nomes, o apagamento dá-se em 90,34\% dos casos. Outro fenômeno que merece ser destacado é o apagamento do gerúndio, que apresenta índices de 96,43\% na variedade, contra 3,57\%, com a preservação do morfema \textit{/ndo/}. Os índices evidenciados nesses dois fenômenos encontram-se acima da média registrada em outras pesquisas discutidas na dissertação. No percurso das pesquisas do mestrado, constatou-se ainda alta produtividade de alteração pronominal, como a de “eles” por \textit{ês}, “ele” por \textit{ê}, “aqueles(as)” por \textit{aquês/aqués}, “você(s)” por \textit{cê(s)} ou \textit{ocê(s)}, bem como de metaplasmos que geram produção de variedades a partir da lógica da redução de palavras, com diferentes tipos de apagamentos. Pode-se concluir que tais fenômenos não são \textit{exclusivos} do território do Vale do Jequitinhonha, mas \textit{também} do Vale, conferindo, assim, características peculiares à dinâmica da língua.
    
Identificar os fenômenos de variação linguística no Vale, por um lado, favorece seu (re)conhecimento e valorização; por outro, indica a necessidade de que esses fenômenos sejam melhor compreendidos, com o intuito de minimizar situações de preconceito linguístico. Nesse sentido, é indispensável considerar a escola como parceira. Para isso, a formação docente deve incluir o estudo de fenômenos linguísticos, de modo que a língua seja compreendida em seu contexto de uso. Torna-se necessário, ainda, que materiais didáticos sejam construídos considerando-se que a língua varia.

\section{Tecnologia digital: Site “Narrativas em Rede”}

A tecnologia digital – o site “Narrativas em Rede” – surge dentro de um contexto mais amplo da dissertação de mestrado intitulada “Variação Linguística em Narrativas Orais no Vale do Jequitinhonha – MG”, no programa de Pós-Graduação em Educação e Docência da FaE/UFMG. Para obtenção do título de mestre, esse programa exige, além da dissertação, a produção de um recurso educacional. Nesse sentido, foi elaborado, no âmbito da pesquisa, o recurso educacional denominado “Narrativas em Rede”. Trata-se de um site cujo objetivo é salvaguardar, da perda e do apagamento histórico, o \textit{corpus} de narrativas orais que foi gravado entre 1987 e 1997, em fitas cassete.

Constatou-se, no decorrer da pesquisa, que, do total das 255 narrativas que foram gravadas pelos projetos, apenas 49 estavam em formato MP3. Dentre elas, citam-se as 35 narrativas disponíveis, em formato MP3, no audiolivro \textit{Negros pelo Vale}, do trabalho de \textcite{Souza2014}, e as 14 faixas dos dois CDs físicos referenciados acima. As demais 206 narrativas encontravam-se em fitas cassete, tecnologia analógica atualmente obsoleta e em desuso.

Objetivando a salvaguarda desse rico patrimônio oral, iniciou-se, em 2024, a digitalização de todo o acervo, cuja meta foi converter todas as narrativas para o formato MP3, criar um site para sitiar o conteúdo e disponibilizar as narrativas (áudios e transcrições), possibilitando usos educativos, estéticos e, artísticos, bem como a preservação do dialetal do Vale do Jequitinhonha (MG), entendido como patrimônio imaterial que precisa preservado e estudado.

Para a extração dos áudios e sua conversão para o formato MP3, foram necessárias a atenção cuidadosa e a escuta minuciosa, a fim de demarcar e localizar com precisão as narrativas, tendo em vista que o conjunto delas se encontrava disperso em 57 fitas cassetes, somando quase uma centena de horas de registros. Para tanto, utilizou-se de dois recursos: (1) um aparelho de som tocador de fitas cassete e (2) o software livre Audacity, instalado em computador, para salvar, cortar e tratar as narrativas. Após centenas de horas de trabalho, pesquisa e dedicação, foram localizadas, no acervo de fitas cassete, 255 narrativas orais, cujas faixas, atualmente em formato MP3, variam de 1 a 44 minutos, de acordo com a estilística de cada narrador.

A narrativa “O Minino Sarvado Ligero”, contada por Silvânio, em Rubim (MG), em 1989, tem, por exemplo, duração de 44 minutos e 15 segundos. Atualmente, todas as 255 narrativas em formato digital MP3 e as 200 transcrições compõem o recurso educacional “Narrativas em Rede”.	Acredita-se que, a partir desse recurso educacional, professores, educadores e escolas possam utilizar as narrativas em áudio e suas transcrições para desenvolver, em consonância com a Base Nacional Comum Curricular \cite{BrasilMEC2018}, as habilidades necessárias à compreensão de que existem diversas variedades linguísticas no português brasileiro e de que todas devem ser valorizadas.

Além disso, o recurso digital propõe evidenciar a tradição oral, o(a) narrador(a) do Vale do Jequitinhonha e, ao mesmo tempo, registrar a oralidade, o dialetal e a estilística típica deste território, pois entende-se que a língua em uso no contexto de 1987 a 1997 apresenta singularidades que pertencem àquelas comunidades. Assim, nas narrativas, entreveem-se memórias do território e, no escopo delas, tais memórias emergem nos modos de narrar e atravessam as histórias de vida do narrador.

Em “Narrativas em Rede”, é possível conhecer um pouco da história do próprio Vale, retratado pela fala e construído e tecido pela oralidade. Como já mencionado, o Vale do Jequitinhonha, com suas tradições e histórias, pode ser compreendido e interpretado pela voz do(a) narrador(a) que viu, viveu e testemunhou fatos e momentos de tempos passados, conforme prospectado por \textcite{Benjamin1985}.

Destaca-se que a criação do site, idealizado 28 anos após as últimas gravações, realizadas no ano de 1997, transcende uma mera iniciativa de digitalização. Trata-se de uma ação de salvaguarda urgente e um gesto de resgate cultural, fundamentado no risco iminente da perda física do acervo. Nesse sentido, a tecnologia digital configura-se como uma janela do presente para o passado e vice-versa.	

O suporte físico original do acervo – as fitas cassete – constitui-se, hoje, uma tecnologia obsoleta, não mais produzida, cujos equipamentos de reprodução se tornam cada vez mais raros e deteriorados. O risco de degradação irreversível do material magnético e analógico aumenta a cada ano, podendo levar ao silenciamento definitivo dessas vozes. A digitalização foi, portanto, uma corrida contra o tempo para evitar a perda total de um patrimônio imaterial único.

O \textit{corpus} de narrativas – composto pelos áudios originais e suas transcrições – constitui-se em um documento da variedade dialetal da região. Ele captura não apenas palavras, mas também entonações, ritmos, sotaques e estruturas gramaticais específicas, que estão em transformação. Preservar os áudios e a estilística típica desses narradores é de suma importância para futuros estudos linguísticos, permitindo compreender a evolução do português brasileiro em uma de suas manifestações mais autênticas, atualmente ameaçadas pela homogeneização midiática.

Salienta-se que, muito antes da escrita se consolidar como modo de transmissão de saber, a oralidade foi, por séculos, o principal mecanismo de circulação de conhecimento, de valores e da própria história. Esse acervo permite, ainda que parcialmente, reviver essa experiência. O site não se configura apenas um repositório, mas uma ponte acústica que permite o contato direto com a oralidade, hoje ameaçada pelos modos digitais e textuais de comunicação. Trata-se, portanto, de uma forma de valorizar e manter viva uma tradição humana fundamental.

Enquanto o acervo físico está confinado entre quatro paredes e sujeito a consultas restritas, o recurso digital “Narrativas em Rede” rompe essas barreiras. Ele transforma um patrimônio local em um bem universalmente acessível, permitindo que pesquisadores, estudantes, descendentes da comunidade e o público em geral, de qualquer parte do mundo, possam acessar, estudar e se emocionar com essas narrativas.

Em um mundo em constante aceleração, a preservação da memória torna-se um ato de resistência. Garantir que o conhecimento, as histórias e as formas de falar do Vale do Jequitinhonha, tal como registradas em 1997, não se percam, é fortalecer a identidade cultural da região. Trata-se de compreender que a inovação e o futuro só se constroem plenamente quando assentados sobre o sólido alicerce do passado.

Portanto, a criação do site “Narrativas em Rede” vai muito além da simples digitalização de um arquivo. É um ato de preservação urgente diante da obsolescência tecnológica. É uma ferramenta vital para a linguística, um testemunho de um modo de transmissão de saber ancestral e um instrumento de democratização do acesso à cultura popular do território do Jequitinhonha. É, em sua essência, garantir que as vozes capturadas há quase três décadas continuem a ecoar, inspirar e ensinar as gerações presentes e futuras.

\printbibliography\label{sec-bib}
% if the text is not in Portuguese, it might be necessary to use the code below instead to print the correct ABNT abbreviations [s.n.], [s.l.]
%\begin{portuguese}
%\printbibliography[title={Bibliography}]
%\end{portuguese}


\begin{dataavailability}
\txtdataavailability{databody} % options: dataavailable, dataonly, databody, datanotav, nodata
\end{dataavailability}


\end{document}

